\documentclass[11pt]{article}

% --------------------------------------------------
% Sprache, Encoding, Typografie
% --------------------------------------------------
\usepackage[ngerman]{babel}
\usepackage[T1]{fontenc}
\usepackage[utf8]{inputenc}
\usepackage{lmodern}
\usepackage{microtype}

% --------------------------------------------------
% Mathematik
% --------------------------------------------------
\usepackage{amsmath,amssymb,amsthm}

% Theorem-Umgebungen (deutsch)
\newtheorem{satz}{Satz}
\newtheorem{lemma}{Lemma}
\newtheorem{korollar}{Korollar}
\theoremstyle{definition}
\newtheorem{definition}{Definition}
\theoremstyle{remark}
\newtheorem{bemerkung}{Bemerkung}
\newtheorem{beispiel}{Beispiel}

% --------------------------------------------------
% Layout
% --------------------------------------------------
\usepackage{geometry}
\geometry{margin=2.5cm}

% --------------------------------------------------
% Links
% --------------------------------------------------
\usepackage{hyperref}

% --------------------------------------------------
% Tabellen
% --------------------------------------------------
\usepackage{booktabs}

% --------------------------------------------------
% Grafiken (robust)
% --------------------------------------------------
\usepackage{graphicx}
\DeclareGraphicsExtensions{.pdf,.png,.jpg}
\graphicspath{{./}{fig/}{figs/}{images/}{img/}}
\newcommand{\includegraphicssafe}[2][]{%
	\IfFileExists{#2}{\includegraphics[#1]{#2}}{%
		\fbox{\parbox[c][0.3\textheight][c]{0.8\textwidth}{\centering Datei nicht gefunden: #2}}%
	}%
}

% --------------------------------------------------
% Titel
% --------------------------------------------------
\title{Primfaktorbasierte Dynamik natürlicher Zahlen\
	\large Eine empirische Analyse der EABC-Struktur}
\author{Thomas Hoffbauer}
\date{\today}

% --------------------------------------------------
\begin{document}
	\maketitle
	
	\begin{abstract}
		Wir untersuchen eine Klassifikation natürlicher Zahlen basierend auf der Zerlegung ihrer Primfaktoren in vier Restklassen modulo 12. Daraus definieren wir eine Signatur $\sigma(n) = (e,a,b,c)$, welche die Häufigkeit der Primfaktoren in den Klassen $E \equiv 1$, $A \equiv 5$, $B \equiv 7$ und $C \equiv 11 \pmod{12}$ beschreibt.
		
		Numerische Experimente zeigen, dass dynamische Eigenschaften arithmetischer Funktionen — insbesondere spektrale Modulation (FFT) und Persistenz (lag1-Autokorrelation) — systematisch durch diese Signatur strukturiert werden.
		
		Dabei tritt eine klare Dichotomie auf: Primfaktoren der Klassen $A$ und $E$ korrelieren mit erhöhter Modulation, während $B$ und $C$ stabilisierend wirken. Die Ergebnisse legen nahe, dass der Zahlraum eine nichttriviale, diskrete Dynamikstruktur besitzt.
	\end{abstract}
	
	% --------------------------------------------------
	\section{Einleitung}
	
	Die klassische Zahlentheorie beschreibt Primzahlen über globale analytische Werkzeuge wie die Riemannsche Zetafunktion. Diese Methoden erfassen jedoch primär globale Verteilungen.
	
	In dieser Arbeit betrachten wir einen komplementären Zugang: die Untersuchung lokaler Strukturen über die Primfaktorzerlegung einzelner Zahlen.
	
	% --------------------------------------------------
	\section{EABC-Signatur}
	
	\begin{definition}
		Für $n \in \mathbb{N}$ mit Primfaktorzerlegung
		\[
		n = \prod_i p_i,
		\]
		definieren wir die Signatur
		\[
		\sigma(n) = (e,a,b,c),
		\]
		wobei
		\[
		\begin{aligned}
			e &= \#\{p_i \equiv 1 \pmod{12}\}, \\
			a &= \#\{p_i \equiv 5 \pmod{12}\}, \\
			b &= \#\{p_i \equiv 7 \pmod{12}\}, \\
			c &= \#\{p_i \equiv 11 \pmod{12}\}.
		\end{aligned}
		\]
	\end{definition}
	
	\begin{bemerkung}
		Diese Klassifikation entspricht der natürlichen Zerlegung der Primzahlen modulo 12 für $p>3$.
	\end{bemerkung}
	
	% --------------------------------------------------
	\section{Dynamische Observablen}
	
	Wir betrachten zwei Größen:
	
	\begin{itemize}
		\item Spektrale Modulation (FFT)
		\item Persistenz (lag1-Autokorrelation)
	\end{itemize}
	
	\begin{definition}
		Die spektrale Modulation ist die dominante Frequenzkomponente einer diskreten Fourier-Transformation.
	\end{definition}
	
	\begin{definition}
		Die Persistenz ist die Autokorrelation erster Ordnung.
	\end{definition}
	
	% --------------------------------------------------
	\section{Empirische Ergebnisse}
	
	\begin{satz}[empirisch]
		Die Größen FFT und lag1 hängen systematisch von der Signatur $\sigma(n)$ ab.
	\end{satz}
	
	\begin{bemerkung}
		Insbesondere gilt:
		\begin{itemize}
			\item Kombinationen von $A$ und $E$ zeigen maximale Modulation.
			\item Dominanz von $B$ und $C$ führt zu hoher Stabilität.
		\end{itemize}
	\end{bemerkung}
	
	% --------------------------------------------------
	\section{Modell}
	
	Wir schlagen folgendes Modell vor:
	\[
	D(n) =
	\alpha_A a + \alpha_E e - \beta_B b - \beta_C c
	+ \gamma_{AE} ae + \gamma_{AC} ac + \gamma_{BC} bc
	\]
	
	\begin{satz}[empirisch]
		Dieses Modell erreicht in numerischen Tests eine Vorhersagegüte von
		\[
		R^2 \approx 0.3 - 0.4.
		\]
	\end{satz}
	
	% --------------------------------------------------
	\section{Schalenmodell}
	
	Die Daten legen eine dreistufige Struktur nahe:
	
	\begin{itemize}
		\item Glatte Zone (hohe Stabilität)
		\item Übergangszone (hohe Modulation)
		\item BC-Kern (strukturstabil)
	\end{itemize}
	
	% --------------------------------------------------
	\section{Diskussion}
	
	Die Ergebnisse zeigen:
	
	\begin{itemize}
		\item Primfaktorstruktur beeinflusst Dynamik messbar
		\item Familien sind nicht äquivalent
		\item Kombinationen erzeugen emergente Effekte
	\end{itemize}
	
	\begin{bemerkung}
		Eine direkte Verbindung zur Zetafunktion wird hier nicht behauptet.
	\end{bemerkung}
	
	% --------------------------------------------------
	\section{Fazit}
	
	\[
	\boxed{
		\text{Der Zahlraum besitzt eine messbare primfaktorbasierte Dynamikstruktur}
	}
	\]
	
	% --------------------------------------------------
	\appendix
	
	\section{Appendix: Beispielcode}
	
	\begin{verbatim}
		import numpy as np
		
		def build_features(sig):
		e,a,b,c = sig
		return np.array([e,a,b,c,ae,ac,b*c])
	\end{verbatim}
	
\end{document}